\subsection{Polynomial Modules from Operators}

Throughout this subsection, $\mathcal{V}$ is finite-dimensional over $\mathbb{F}$ unless explicitly stated otherwise.

\begin{definition} \label{def:operator-module-and-matrix-module} \hlt{The Module Attached to an Operator}\\
For $T\in\LL(\mathcal{V})$, the \hlt{operator module} attached to $T$ is the $\mathbb{F}[z]$-module $\mathcal{V}_T$ whose underlying $\mathbb{F}$-vector space is $\mathcal{V}$.\\
Its scalar action is $p(z)\cdot\mathbf{v}\defeq p(T)\mathbf{v}$. For $A\in\operatorname{Mat}(n,n,\mathbb{F})$, define $(\mathbb{F}^n)_A$ by $p(z)\cdot\mathbf{x}\defeq p(A)\mathbf{x}$.
\end{definition}

\begin{proposition} \label{prop:operator-modules-are-fg-torsion} \hlt{Operator Modules Are Finitely Generated Torsion Modules}\\
Let $0<\dim\mathcal{V}<\infty$ and $T\in\LL(\mathcal{V})$.
\begin{enumerate}[label=(\roman*)]
\setlength{\itemsep}{0pt}
\item $\mathcal{V}_T$ is finitely generated over $\mathbb{F}[z]$.
\item $\mathcal{V}_T$ is a torsion $\mathbb{F}[z]$-module.
\end{enumerate}
\begin{proof}
\begin{enumerate}[label=(\roman*)]
\setlength{\itemsep}{0pt}
\item Let $(\mathbf{v}_1,\ldots,\mathbf{v}_n)$ be an $\mathbb{F}$-basis of $\mathcal{V}$. Then $\mathcal{V}_T=\langle\!\langle\mathbf{v}_1,\ldots,\mathbf{v}_n\rangle\!\rangle_{\mathbb{F}[z]}$ by $\mathbb{F}\subseteq\mathbb{F}[z]$.
\item Choose $0\neq p\in\mathbb{F}[z]$ with $p(T)=0$ by Proposition \ref{prop:finite-dimensional-operators-have-annihilating-polynomials}; $(\forall\,\mathbf{v}\in\mathcal{V})\,p\cdot\mathbf{v}=p(T)\mathbf{v}=\mathbf{0}$, so $\mathcal{V}_T$ is torsion by Definition \ref{def:module-torsion}(iii).
\end{enumerate}
\hlt{Structural Commentary}: An $\mathbb{F}$-basis handles generation; one operator annihilator handles torsion.\\
\hlt{Key Insight}: One nonzero polynomial annihilates every vector.
\end{proof}
\end{proposition}

\begin{theorem} \label{thm:operator-submodules-invariant-subspaces} \hlt{Submodules Are Invariant Subspaces}\\
Let $T\in\LL(\mathcal{V})$ and $\mathcal{S}\subseteq\mathcal{V}$. Then $\mathcal{S}\leq_{\mathbb{F}[z]}\mathcal{V}_T\Longleftrightarrow(\mathcal{S}\leq_{\mathbb{F}}\mathcal{V}\land T\mathcal{S}\subseteq\mathcal{S})$.
\begin{proof}
\begin{enumerate}[label=\roman*.]
\setlength{\itemsep}{0pt}
\item[$(\Rightarrow)$] Suppose $\mathcal{S}\leq_{\mathbb{F}[z]}\mathcal{V}_T$. Then $\mathcal{S}\leq_{\mathbb{F}}\mathcal{V}$ by $\mathbb{F}\subseteq\mathbb{F}[z]$. Hence $T\mathcal{S}\subseteq\mathcal{S}$ by $(\forall\,\mathbf{s}\in\mathcal{S})\,T\mathbf{s}=z\cdot\mathbf{s}\in\mathcal{S}$.
\item[$(\Leftarrow)$] Suppose $\mathcal{S}\leq_{\mathbb{F}}\mathcal{V}$ and $T\mathcal{S}\subseteq\mathcal{S}$. Then $(\forall\,j\in\mathbb{Z}_{\geq0})\,T^j\mathcal{S}\subseteq\mathcal{S}$ by induction on $j$.\\
For $p(z)=\sum\limits_{j=0}^{m}a_jz^j\in\mathbb{F}[z]$ and $\mathbf{s}\in\mathcal{S}$, $p\cdot\mathbf{s}=\sum\limits_{j=0}^{m}a_jT^j\mathbf{s}\in\mathcal{S}$. Thus $\mathcal{S}\leq_{\mathbb{F}[z]}\mathcal{V}_T$.
\end{enumerate}
\hlt{Structural Commentary}: Restriction to constants recovers linearity; closure under $z$ extends to every polynomial.\\
\hlt{Key Insight}: The scalar $z$ acts exactly as $T$.
\end{proof}
\end{theorem}

\begin{theorem} \label{thm:operator-similarity-iff-module-isomorphism} \hlt{Similarity Is Module Isomorphism}\\
Let $T,S\in\LL(\mathcal{V})$.
\begin{enumerate}[label=(\roman*)]
\setlength{\itemsep}{0pt}
\item For $\Phi\in\operatorname{GL}(\mathcal{V})$, $\Phi\in\operatorname{Hom}_{\mathbb{F}[z]}(\mathcal{V}_T,\mathcal{V}_S)\Longleftrightarrow\Phi T=S\Phi$.
\item $\mathcal{V}_T\cong_{\mathbb{F}[z]}\mathcal{V}_S\Longleftrightarrow T\sim S$.
\end{enumerate}
\begin{proof}
\begin{enumerate}[label=(\roman*)]
\setlength{\itemsep}{0pt}
\item[$(i)$] ($\Rightarrow$) If $\Phi:\mathcal{V}_T\to\mathcal{V}_S$ is $\mathbb{F}[z]$-linear, then $(\forall\,\mathbf{v}\in\mathcal{V})\,\Phi T\mathbf{v}=\Phi(z\cdot\mathbf{v})=z\cdot\Phi\mathbf{v}=S\Phi\mathbf{v}$. Hence $\Phi T=S\Phi$.
\item[$(i)$] ($\Leftarrow$) Suppose $\Phi T=S\Phi$. Then $(\forall\,j\in\mathbb{Z}_{\geq0})\,\Phi T^j=S^j\Phi$ by induction on $j$.\\
Thus $(\forall\,(p,\mathbf{v})\in\mathbb{F}[z]\times\mathcal{V})\,\Phi(p\cdot\mathbf{v})=\Phi(p(T)\mathbf{v})=p(S)\Phi\mathbf{v}=p\cdot\Phi\mathbf{v}$, so $\Phi$ is $\mathbb{F}[z]$-linear.
\item[$(ii)$] ($\Rightarrow$) Let $\Phi:\mathcal{V}_T\xrightarrow{\sim}\mathcal{V}_S$ be an $\mathbb{F}[z]$-module isomorphism. Then $\Phi\in\operatorname{GL}(\mathcal{V})$ by $\mathbb{F}\subseteq\mathbb{F}[z]$, and $\Phi T=S\Phi$ by part (i).\\
Hence $S=\Phi T\Phi^{-1}$, so $T\sim S$ by Definition \ref{def:similar-linear-operators}.
\item[$(ii)$] ($\Leftarrow$) Choose $\Phi\in\operatorname{GL}(\mathcal{V})$ with $S=\Phi T\Phi^{-1}$ by Definition \ref{def:similar-linear-operators}.\\
Then $\Phi T=S\Phi$, so $\Phi$ is $\mathbb{F}[z]$-linear by part (i). Since $\Phi$ is bijective, $\Phi:\mathcal{V}_T\xrightarrow{\sim}\mathcal{V}_S$.
\end{enumerate}
\hlt{Structural Commentary}: Intertwining the action of $z$ extends to every polynomial and characterizes conjugacy.\\
\hlt{Key Insight}: Module linearity is exactly the equation $\Phi T=S\Phi$.
\end{proof}
\end{theorem}

\begin{corollary} \label{cor:matrix-similarity-iff-matrix-modules-isomorphic} \hlt{Matrix Similarity Is Matrix-Module Isomorphism}\\
Let $n\geq1$ and $A,B\in\operatorname{Mat}(n,n,\mathbb{F})$.
\begin{enumerate}[label=(\roman*)]
\setlength{\itemsep}{0pt}
\item For $X\in\operatorname{GL}_n(\mathbb{F})$, $T_X\in\operatorname{Hom}_{\mathbb{F}[z]}((\mathbb{F}^n)_A,(\mathbb{F}^n)_B)\Longleftrightarrow XA=BX$.
\item $(\mathbb{F}^n)_A\cong_{\mathbb{F}[z]}(\mathbb{F}^n)_B\Longleftrightarrow A\sim B$.
\end{enumerate}
\begin{proof}
\begin{enumerate}[label=(\roman*)]
\setlength{\itemsep}{0pt}
\item[$(i)$] ($\Rightarrow$) If $T_X$ is $\mathbb{F}[z]$-linear, then $T_XT_A=T_BT_X$ by Theorem \ref{thm:operator-similarity-iff-module-isomorphism}(i). Hence $(\forall\,\mathbf{x}\in\mathbb{F}^n)\,XA\mathbf{x}=BX\mathbf{x}$, so $XA=BX$.
\item[$(i)$] ($\Leftarrow$) If $XA=BX$, then $T_XT_A=T_BT_X$ by Definition \ref{def:linear-map-induced-by-matrix}. Thus $T_X$ is $\mathbb{F}[z]$-linear by Theorem \ref{thm:operator-similarity-iff-module-isomorphism}(i).
\item[$(ii)$] ($\Rightarrow$) Let $\Psi:(\mathbb{F}^n)_A\xrightarrow{\sim}(\mathbb{F}^n)_B$. Choose $X,Y\in\operatorname{Mat}(n,n,\mathbb{F})$ with $\Psi=T_X$ and $\Psi^{-1}=T_Y$ by Theorem \ref{thm:linear-maps-and-matrices-same-data}(i).\\
For $\mathbf{x}\in\mathbb{F}^n$, $XY\mathbf{x}=YX\mathbf{x}=\mathbf{x}$ by the inverse identities and Definition \ref{def:linear-map-induced-by-matrix}. Hence $XY=YX=I_n$ and $X\in\operatorname{GL}_n(\mathbb{F})$.\\
Then $XA=BX$ by part (i), and $A\sim B$ by $B=XAX^{-1}$ and Definition \ref{def:similar-matrices}.
\item[$(ii)$] ($\Leftarrow$) Choose $X\in\operatorname{GL}_n(\mathbb{F})$ with $B=XAX^{-1}$ by Definition \ref{def:similar-matrices}.\\
Then $XA=BX$, so the bijection $T_X$ is $\mathbb{F}[z]$-linear by part (i) and is an $\mathbb{F}[z]$-module isomorphism.
\end{enumerate}
\hlt{Structural Commentary}: Coordinate intertwiners translate module isomorphisms into the conjugation equation $XA=BX$.\\
\hlt{Key Insight}: Matrix similarity is polynomial-module isomorphism in coordinates.
\end{proof}
\end{corollary}

\begin{definition} \label{def:operator-module-annihilator-minimal-polynomial} \hlt{Operator-Module Annihilator}\\
Let $0<\dim\mathcal{V}<\infty$ and $T\in\LL(\mathcal{V})$. The \hlt{annihilator} of $\mathcal{V}_T$ is $\operatorname{Ann}_{\mathbb{F}[z]}(\mathcal{V}_T)\defeq\{p\in\mathbb{F}[z]:p(T)=0_{\LL(\mathcal{V})}\}$.
\end{definition}

\begin{proposition} \label{prop:minimal-polynomial-module-order} \hlt{Minimal Polynomial as Module Order}\\
Let $0<\dim\mathcal{V}<\infty$ and $T\in\LL(\mathcal{V})$.
\begin{enumerate}[label=(\roman*)]
\setlength{\itemsep}{0pt}
\item $\operatorname{Ann}_{\mathbb{F}[z]}(\mathcal{V}_T)=\langle\mu_T\rangle$. Equivalently, $\mu_T$ is unique monic polynomial of minimal degree in $\{p\in\mathbb{F}[z]:p(T)=0\}$.
\item If $S\in\LL(\mathcal{V})$ and $S\sim T$, then $\mu_S=\mu_T$.
\item If $n\geq1$ and $A,B\in\operatorname{Mat}(n,n,\mathbb{F})$ are similar, then $\mu_A=\mu_B$.
\end{enumerate}
\begin{proof}
\begin{enumerate}[label=(\roman*)]
\setlength{\itemsep}{0pt}
\item $\operatorname{Ann}_{\mathbb{F}[z]}(\mathcal{V}_T)=\{p\in\mathbb{F}[z]:p(T)=0\}$ by Definition \ref{def:operator-module-annihilator-minimal-polynomial}.\\
For $p\neq0$, $p(T)=0\Longleftrightarrow\mu_T\mid p$ by Theorem \ref{thm:operator-annihilator-iff-multiple-minimal-polynomial}; the zero polynomial belongs to both ideals.\\
Hence $\operatorname{Ann}_{\mathbb{F}[z]}(\mathcal{V}_T)=\langle\mu_T\rangle$, and the characterization of $\mu_T$ follows by Definition \ref{def:operator-minimal-polynomial}.
\item Choose $\Phi\in\operatorname{GL}(\mathcal{V})$ with $S=\Phi T\Phi^{-1}$ by Definition \ref{def:similar-linear-operators}. Then $(\forall\,p\in\mathbb{F}[z])\,p(S)=\Phi p(T)\Phi^{-1}$, so $\operatorname{Ann}_{\mathbb{F}[z]}(\mathcal{V}_S)=\operatorname{Ann}_{\mathbb{F}[z]}(\mathcal{V}_T)$.\\
Hence $\mu_S=\mu_T$ by part (i).
\item $\mu_A=\mu_B$ by Theorem \ref{thm:similarity-preserves-minimal-polynomial}(ii).
\end{enumerate}
\hlt{Structural Commentary}: Annihilator divisibility identifies $\mu_T$; conjugation transports the resulting principal ideal.\\
\hlt{Key Insight}: The module order is invariant under intertwiners.
\end{proof}
\end{proposition}

\begin{definition} \label{def:cyclic-subspace-and-vector-annihilator} \hlt{Cyclic Subspace and Vector Annihilator}\\
Let $T\in\LL(\mathcal{V})$ and $\mathbf{v}\in\mathcal{V}$.
\begin{enumerate}[label=(\roman*)]
\setlength{\itemsep}{0pt}
\item The \hlt{$T$-cyclic subspace generated by $\mathbf{v}$} is $\mathcal{C}_T(\mathbf{v})\defeq\operatorname{span}\{T^j\mathbf{v}:j\in\mathbb{Z}_{\geq0}\}$.
\item The vector $\mathbf{v}$ is a \hlt{cyclic vector} for $T$ if $\mathcal{C}_T(\mathbf{v})=\mathcal{V}$.
\item The \hlt{annihilator of $\mathbf{v}$ for $T$} is $\operatorname{ann}_T(\mathbf{v})\defeq\{p\in\mathbb{F}[z]:p(T)\mathbf{v}=\mathbf{0}\}$.
\end{enumerate}
\end{definition}

\begin{proposition} \label{prop:cyclic-subspace-smallest-invariant-subspace} \hlt{Cyclic Subspaces Are Smallest Invariant Hosts}\\
Let $T\in\LL(\mathcal{V})$ and $\mathbf{v}\in\mathcal{V}$.
\begin{enumerate}[label=(\roman*)]
\setlength{\itemsep}{0pt}
\item $\mathcal{C}_T(\mathbf{v})=\{p(T)\mathbf{v}:p\in\mathbb{F}[z]\}$.
\item $\mathcal{C}_T(\mathbf{v})$ is the smallest $T$-invariant subspace containing $\mathbf{v}$.
\item $\operatorname{ann}_T(\mathbf{v})\trianglelefteq\mathbb{F}[z]$.
\end{enumerate}
\begin{proof}
\begin{enumerate}[label=(\roman*)]
\setlength{\itemsep}{0pt}
\item For $p(z)=\sum\limits_{j=0}^{m}a_jz^j\in\mathbb{F}[z]$, $p(T)\mathbf{v}=\sum\limits_{j=0}^{m}a_jT^j\mathbf{v}$. Both inclusions hold by this identity.
\item $T(p(T)\mathbf{v})=(zp)(T)\mathbf{v}\in\mathcal{C}_T(\mathbf{v})$ by part (i), so $\mathcal{C}_T(\mathbf{v})$ is $T$-invariant.\\
Let $\mathcal{U}\leq_{\mathbb{F}}\mathcal{V}$ contain $\mathbf{v}$ and satisfy $T\mathcal{U}\subseteq\mathcal{U}$. By induction, $(\forall\,j\in\mathbb{Z}_{\geq0})\,T^j\mathbf{v}\in\mathcal{U}$, so $\mathcal{C}_T(\mathbf{v})\subseteq\mathcal{U}$.
\item The zero polynomial belongs to $\operatorname{ann}_T(\mathbf{v})$.\\
For $a,b\in\mathbb{F}[z]$ and $p,q\in\operatorname{ann}_T(\mathbf{v})$, $(ap+bq)(T)\mathbf{v}=a(T)p(T)\mathbf{v}+b(T)q(T)\mathbf{v}=\mathbf{0}$ by Proposition \ref{prop:operator-polynomial-products}(i).\\
Thus $\operatorname{ann}_T(\mathbf{v})\leq{}_{\mathbb{F}[z]}\mathbb{F}[z]$ by Theorem \ref{thm:submodule-test}, and $\operatorname{ann}_T(\mathbf{v})\trianglelefteq\mathbb{F}[z]$ by Proposition \ref{prop:submodules-of-regular-module-are-ideals}.
\end{enumerate}
\hlt{Structural Commentary}: Orbit-polynomial equality controls minimal invariance; evaluation closure controls the ideal property.\\
\hlt{Key Insight}: Polynomial evaluation packages every finite orbit combination.
\end{proof}
\end{proposition}

\begin{perspective} \label{per:operator-as-polynomial-action} \hlt{An Operator as a Polynomial Action}\\
Definition \ref{def:operator-module-and-matrix-module} interprets $z$ as $T$. Proposition \ref{prop:cyclic-subspace-smallest-invariant-subspace}(i) identifies the cyclic module generated by $\mathbf{v}$ with $\mathcal{C}_T(\mathbf{v})$.\\
Figure \ref{fig:operator-cyclic-subspace-polynomial-action} illustrates how iteration under $T$ forms this cyclic span; the cited formal owners establish the identities.
\end{perspective}

\begin{illustrationfigure}[H]
  \centering
  \begin{tikzpicture}[
  >=Stealth,
  every node/.style={font=\small},
  orbit/.style={draw=blue!60!black, fill=blue!6, rounded corners=2pt,
    minimum width=1.15cm, minimum height=0.58cm, align=center, inner sep=4pt},
  subspace/.style={draw=dartmouthgreen!65!black, fill=dartmouthgreen!8, rounded corners=2pt,
    minimum width=4.7cm, minimum height=0.72cm, align=center, inner sep=6pt},
  arr/.style={->, line width=0.7pt, draw=black!70}
]
  \node[orbit] (v0) at (-2.45,0) {$\mathbf{v}$};
  \node[orbit] (v1) at (-0.85,0) {$T\mathbf{v}$};
  \node[orbit] (v2) at (0.85,0) {$T^2\mathbf{v}$};
  \node[orbit, minimum width=0.8cm] (dots) at (2.35,0) {$\cdots$};

  \draw[arr] (v0) -- node[above] {$T$} (v1);
  \draw[arr] (v1) -- node[above] {$T$} (v2);
  \draw[arr] (v2) -- node[above] {$T$} (dots);

  \node[subspace] (C) at (0,-1.35)
    {$\mathcal{C}_T(\mathbf{v})=\operatorname{span}\{T^j\mathbf{v}:j\geq0\}$};

  \draw[arr] (0,-0.38) -- node[right] {$\operatorname{span}$} (C.north);
\end{tikzpicture}

  \caption{Iteration under $T$ forms the orbit whose finite linear span is $\mathcal{C}_T(\mathbf{v})$.}
  \label{fig:operator-cyclic-subspace-polynomial-action}
\end{illustrationfigure}

\begin{definition} \label{def:vector-order-under-operator} \hlt{Order of a Vector under an Operator}\\
Let $T\in\LL(\mathcal{V})$ and $\mathbf{v}\in\mathcal{V}\setminus\{\mathbf{0}\}$. Then $\{0\}\neq\operatorname{ann}_T(\mathbf{v})\trianglelefteq\mathbb{F}[z]$ by Proposition \ref{prop:operator-modules-are-fg-torsion}(ii), Definition \ref{def:module-torsion}(iii), Proposition \ref{prop:cyclic-subspace-smallest-invariant-subspace}(iii).\\
This ideal has a unique monic generator by Theorem \ref{thm:integers-and-polynomial-ring-are-pids}(iii)--(iv).\\
Specializing Definition \ref{def:module-order-over-pid}, the \hlt{order of $\mathbf{v}$ under $T$} is this generator $m_{\mathbf{v}}$; equivalently, $\operatorname{ann}_T(\mathbf{v})=\langle m_{\mathbf{v}}\rangle$.
\end{definition}

\begin{definition} \label{def:companion-matrix} \hlt{Companion Matrix}\\
Let $n\geq1$ and let $f(z)=z^n+c_{n-1}z^{n-1}+\cdots+c_1z+c_0$ be monic. For $n\geq2$, the \hlt{companion matrix} of $f$ is
\begin{align}
C_f\defeq
\begin{bmatrix}
0&0&\cdots&0&-c_0\\
1&0&\cdots&0&-c_1\\
0&1&\cdots&0&-c_2\\
\vdots&\vdots&\ddots&\vdots&\vdots\\
0&0&\cdots&1&-c_{n-1}
\end{bmatrix}.
\nonumber
\end{align}
For $n=1$, set $C_f=[-c_0]$.
\end{definition}

\begin{theorem} \label{thm:companion-matrix-has-prescribed-minimal-polynomial} \hlt{Companion Matrix Realises a Polynomial}\\
Let $n\geq1$, let $f(z)=z^n+\sum\limits_{i=0}^{n-1}c_iz^i$ be monic, and let $C_f$ be its companion matrix.
\begin{enumerate}[label=(\roman*)]
\setlength{\itemsep}{0pt}
\item $\mu_{C_f}=f$.
\item $p_{C_f}=f$.
\end{enumerate}
\begin{proof}
\begin{enumerate}[label=(\roman*)]
\setlength{\itemsep}{0pt}
\item Let $(\mathbf{e}_1,\ldots,\mathbf{e}_n)$ be the standard basis. Then $(\forall\,j\in\{0,\ldots,n-1\})\,C_f^j\mathbf{e}_1=\mathbf{e}_{j+1}$.\\
The last column satisfies $C_f^n\mathbf{e}_1=-\sum\limits_{i=0}^{n-1}c_iC_f^i\mathbf{e}_1$, so $f(C_f)\mathbf{e}_1=\mathbf{0}$.\\
For $j\in\{1,\ldots,n\}$, $f(C_f)\mathbf{e}_j=C_f^{j-1}f(C_f)\mathbf{e}_1=\mathbf{0}$ by $\mathbf{e}_j=C_f^{j-1}\mathbf{e}_1$ and commutation with $C_f$. Thus $f(C_f)=0$.\\
If $0\neq g(z)=\sum\limits_{i=0}^{m}b_iz^i$ with $m<n$, then $g(C_f)\mathbf{e}_1=\sum\limits_{i=0}^{m}b_i\mathbf{e}_{i+1}\neq\mathbf{0}$ by linear independence of the standard basis.\\
Therefore no nonzero polynomial of degree less than $n$ annihilates $C_f$.\\
Since $f$ is a monic degree-$n$ annihilator, $\mu_{C_f}=f$ by Definition \ref{def:minimal-polynomial-matrix}.
\item Work over the rational-function field $\mathbb{F}(z)$ and set $M(z)\defeq zI_n-C_f$.\\
Row operation $R_1\leftarrow R_1+\sum\limits_{k=2}^{n}z^{k-1}R_k$ preserves $\det M(z)$ by Theorem \ref{thm:determinant-basic-laws}(ii)--(iv),(ix). New first row is $(0,\ldots,0,f(z))$.\\
Moving it below the other rows uses $n-1$ swaps and produces an upper-triangular matrix with diagonal $(-1,\ldots,-1,f(z))$.\\
Hence $(-1)^{n-1}\det M(z)=(-1)^{n-1}f(z)$ by Theorem \ref{thm:determinant-basic-laws}(iii),(v), so $p_{C_f}(z)=\det M(z)=f(z)$ by Definition \ref{def:matrix-characteristic-polynomial}.
\end{enumerate}
\hlt{Structural Commentary}: The orbit basis fixes the least annihilator; determinant-preserving row reduction computes $p_{C_f}$.\\
\hlt{Key Insight}: One last-column relation controls both polynomial invariants.
\end{proof}
\end{theorem}

\begin{theorem} \label{thm:cyclic-subspace-quotient-and-companion} \hlt{Cyclic Subspaces as Polynomial Quotients}\\
Let $\dim\mathcal{V}<\infty$, $T\in\LL(\mathcal{V})$, and $0\neq\mathbf{v}\in\mathcal{V}$. Let $m_{\mathbf{v}}(z)=z^d+\sum\limits_{i=0}^{d-1}c_iz^i$ be the order of $\mathbf{v}$, where $d\geq1$.\\
Define $\Phi_{\mathbf{v}}:\mathbb{F}[z]\to\mathcal{C}_T(\mathbf{v})$ by $\Phi_{\mathbf{v}}(p)=p(T)\mathbf{v}$. Set $\mathcal{B}_{\mathbf{v}}\defeq(\mathbf{v},T\mathbf{v},\ldots,T^{d-1}\mathbf{v})$. Then
\begin{enumerate}[label=(\roman*)]
\setlength{\itemsep}{0pt}
\item $\Phi_{\mathbf{v}}$ is a surjective $\mathbb{F}[z]$-module homomorphism with $\ker\Phi_{\mathbf{v}}=\langle m_{\mathbf{v}}\rangle$, and $\mathbb{F}[z]/\langle m_{\mathbf{v}}\rangle\cong\mathcal{C}_T(\mathbf{v})$;
\item $\mathcal{B}_{\mathbf{v}}$ is a basis of $\mathcal{C}_T(\mathbf{v})$;
\item $[T|_{\mathcal{C}_T(\mathbf{v})}]_{\mathcal{B}_{\mathbf{v}}\leftarrow\mathcal{B}_{\mathbf{v}}}=C_{m_{\mathbf{v}}}$.
\end{enumerate}
\begin{proof}
\begin{enumerate}[label=(\roman*)]
\setlength{\itemsep}{0pt}
\item For $a,b,p,q\in\mathbb{F}[z]$, $\Phi_{\mathbf{v}}(ap+bq)=a\cdot\Phi_{\mathbf{v}}(p)+b\cdot\Phi_{\mathbf{v}}(q)$ by Proposition \ref{prop:operator-polynomial-products}(i).\\
Thus $\Phi_{\mathbf{v}}$ is $\mathbb{F}[z]$-linear by Definition \ref{def:module-homomorphism}(i). Its image is $\mathcal{C}_T(\mathbf{v})$ by Proposition \ref{prop:cyclic-subspace-smallest-invariant-subspace}(i).\\
Its kernel is $\operatorname{ann}_T(\mathbf{v})=\langle m_{\mathbf{v}}\rangle$ by Definition \ref{def:vector-order-under-operator}. The quotient isomorphism follows by Theorem \ref{thm:module-isomorphism-theorems}(i).
\item If $\sum\limits_{j=0}^{d-1}a_jT^j\mathbf{v}=\mathbf{0}$, set $q(z)\defeq\sum\limits_{j=0}^{d-1}a_jz^j\in\langle m_{\mathbf{v}}\rangle$. If $q\neq0$, then $m_{\mathbf{v}}\mid q$ and $d\leq\deg q<d$, a contradiction.\\
Thus $q=0$, so $\mathcal{B}_{\mathbf{v}}$ is independent. For $p\in\mathbb{F}[z]$, choose $h,r\in\mathbb{F}[z]$ with $p=hm_{\mathbf{v}}+r$ and $\deg r<d$ or $r=0$ by Theorem \ref{thm:polynomial-division-algorithm}(i).\\
Then $p(T)\mathbf{v}=r(T)\mathbf{v}\in\operatorname{span}\mathcal{B}_{\mathbf{v}}$, so $\mathcal{B}_{\mathbf{v}}$ spans.
\item For $0\leq j<d-1$, $T(T^j\mathbf{v})=T^{j+1}\mathbf{v}$. The relation $T^d\mathbf{v}=-\sum\limits_{i=0}^{d-1}c_iT^i\mathbf{v}$ holds by $m_{\mathbf{v}}(T)\mathbf{v}=\mathbf{0}$.\\
These are precisely the columns of $C_{m_{\mathbf{v}}}$ by Definitions \ref{def:matrix-representation-of-linear-map} and \ref{def:companion-matrix}.
\end{enumerate}
\hlt{Structural Commentary}: The evaluation kernel forms the quotient; division selects the basis; the order relation fixes the matrix.\\
\hlt{Key Insight}: The first orbit relation becomes the companion column.
\end{proof}
\end{theorem}

\begin{perspective} \label{per:module-operator-dictionary} \hlt{The Translation Dictionary}\\
Definitions \ref{def:operator-module-and-matrix-module}, \ref{def:cyclic-subspace-and-vector-annihilator}, and \ref{def:vector-order-under-operator} identify the module objects.\\
Proposition \ref{prop:minimal-polynomial-module-order} and Theorems \ref{thm:operator-submodules-invariant-subspaces} and \ref{thm:cyclic-subspace-quotient-and-companion} establish the corresponding operator structures:\\
\noindent\begin{tabularx}{\textwidth}{@{}p{0.48\textwidth}|X@{}}
$\mathbb{F}[z]$-module $\mathcal{V}_T$ & $\mathbb{F}$-vector space $\mathcal{V}$\\
\hline
$p(z)\cdot\mathbf{v}$ & $p(T)\mathbf{v}$\\
submodule & $T$-invariant subspace\\
$\operatorname{Ann}_{\mathbb{F}[z]}(\mathcal{V}_T)=\langle\mu_T\rangle$ & minimal polynomial of $T$\\
$\langle\!\langle\mathbf{v}\rangle\!\rangle_{\mathbb{F}[z]}$ & $\mathcal{C}_T(\mathbf{v})$\\
$\mathbb{F}[z]/\langle m_{\mathbf{v}}\rangle$ & cyclic block with companion matrix $C_{m_{\mathbf{v}}}$
\end{tabularx}
This dictionary is local to $\mathcal{V}_T$; classification of arbitrary operator modules requires the subsequent PID decomposition results.
\end{perspective}
